%=======================================================================
%  IESE EXERCISE BOOK 3 -- POWER SUPPLIES AND REGULATORS
%  Problems to work on (no worked method); answers at the back.
%  Companion to iese03_power_manual.tex.
%=======================================================================
\documentclass[11pt,a4paper,oneside]{book}
\usepackage{iese_manual}

\hypersetup{pdftitle={IESE Exercise Book 3 -- Power}}

% answer-key helper
\newcommand{\ans}[1]{\par\smallskip\noindent\textbf{#1}\quad\ignorespaces}

\begin{document}
\frontmatter

\begingroup
\thispagestyle{empty}\centering
\vspace*{2.2cm}
{\color{accent}\rule{\linewidth}{1.4pt}}\\[0.5cm]
{\Huge\bfseries\color{accent} Power Supplies \& Regulators}\\[0.35cm]
{\LARGE Exercise Book}\\[0.4cm]
{\color{accent}\rule{\linewidth}{1.4pt}}\\[1.2cm]
{\large IESE Exercise Book 3 of 3}\\[1.6cm]
{\footnotesize Problems only --- attempt them first; answers (final results, no
worked steps) are at the back.}
\par\vspace*{\fill}
\endgroup
\clearpage

\chapter*{How to use this book}
\addcontentsline{toc}{chapter}{How to use this book}
\markboth{How to use this book}{How to use this book}

This is the problem set distilled from the \emph{Power Supplies \& Regulators}
solving manual. It contains \textbf{only exercises and their answers} --- no
theory and no worked derivations. Each chapter mirrors a chapter of the manual,
so when you are stuck on a \emph{method}, open the manual at the same chapter and
read its Protocol and Worked Example.

\begin{keyidea}
Sizing exercises run \emph{backwards}, from the load to the wall: fix the
regulator's headroom first, then the ripple, then the diode drops, then the
worst-case mains. The answer key gives the standard part/value reached at the
end --- get there yourself before you look.
\end{keyidea}

\tableofcontents
\clearpage
\mainmatter

%=======================================================================
\chapter{Specifications: reading the data sheet}
%=======================================================================

\begin{exercise}[A datasheet line by line]
A \qty{5}{\volt} regulator's table reads: $\Delta V_{out}=\qty{10}{\milli\volt}$
for $I_{out}$ from \qty{5}{\milli\ampere} to \qty{1.5}{\ampere};
$\Delta V_{out}=\qty{3}{\milli\volt}$ for $V_{in}$ from \qtyrange{8}{12}{\volt}.
(a)~Load and line sensitivity, and the ripple rejection in dB.
(b)~At $V_{in}=\qty{12}{\volt}$, $I_{out}=\qty{1}{\ampere}$: efficiency and
dissipated power.
\end{exercise}

\begin{exercise}[Reading the numbers]
A regulator shows $\Delta V_{out}=\qty{7}{\milli\volt}$ over an input range
$\Delta V_{in}=\qty{17.5}{\volt}$, and $\Delta V_{out}=\qty{30}{\milli\volt}$ for
a load step $\Delta I_{out}=\qty{1.5}{\ampere}$. Find the ripple rejection (dB)
and the output resistance.
\end{exercise}

%=======================================================================
\chapter{The classical (transformer) power supply}
%=======================================================================

\begin{exercise}[Size the secondary]
Size the supply for \qty{12}{\volt}/\qty{0.5}{\ampere} from a 7812
($V_{DO,min}=\qty{2}{\volt}$), full-wave bridge with $V_D=\qty{1}{\volt}$, ripple
$\Delta V\le\qty{1}{\volt}$, mains $230\pm\qty{10}{\percent}$ at \qty{50}{\hertz}.
(a)~The filter capacitor.
(b)~The secondary RMS voltage and the turns ratio.
(c)~Why would a \emph{half-wave} rectifier double the capacitor?
\end{exercise}

\begin{exercise}[Filter capacitor]
Full-wave, $f=\qty{50}{\hertz}$, $V_{MAX}=\qty{9}{\volt}$,
$I_{out}=\qty{100}{\milli\ampere}$, $C=\qty{2.2}{\milli\farad}$ (use
$T_{dis}=T$). Find the ripple $\Delta V$, the conduction time $t_c$ and the peak
charging current $I_{peak}$.
\end{exercise}

\begin{exercise}[The capacitor trade-off]
Repeat the previous exercise with $C'=\qty{4.7}{\milli\farad}$ (same
$f,V_{MAX},I_{out}$). What happens to $\Delta V$ and to $I_{peak}$?
\end{exercise}

%=======================================================================
\chapter{Power factor and its correction}
%=======================================================================

\begin{exercise}[$R$--$L$ load and its correction]
$R=\qty{3}{\ohm}$, $L=\qty{10.6}{\milli\henry}$, $V_S=\qty{480}{\volt_{RMS}}$ at
$f=\qty{60}{\hertz}$ ($\omega=\qty{377}{rad/s}$).
(a)~Find $|I_S|$, $\varphi$, $S$, $P$ and the power factor.
(b)~Find the parallel capacitor that corrects $PF$ to $1$, and the new source
current.
\end{exercise}

\begin{exercise}[The pure inductor]
A purely inductive load $L=\qty{10.6}{\milli\henry}$ (no $R$), same source as
above. (a)~Find $|I_S|$ and the power factor.
(b)~Find the parallel $C$ that brings the source current to zero.
\end{exercise}

%=======================================================================
\chapter{Linear voltage regulators}
%=======================================================================

\begin{exercise}[7805 at light load]
A 7805 supplies $I_{out}=\qty{200}{\milli\ampere}$ from $V_{in}=\qty{8}{\volt}$.
(a)~Dissipation, efficiency, and is the bare package ($R_{th}=\qty{60}{K/W}$)
enough at $T_{amb}=\qty{25}{\celsius}$?
(b)~What is the lowest rippled-input \emph{valley} that still regulates?
\end{exercise}

\begin{exercise}[7805 needs a heat sink]
A 7805 supplies $I_{out}=\qty{1}{\ampere}$ from $V_{in}=\qty{12}{\volt}$.
Find the dissipation and efficiency, and the junction temperature rise both bare
($\qty{60}{K/W}$) and with a \qty{10}{K/W} heat sink.
\end{exercise}

\begin{exercise}[Efficiency ceiling]
For a 78xx ($V_{DO,min}=\qty{2}{\volt}$), give the maximum efficiency
$\eta_{MAX}=V_{out}/(V_{out}+2)$ for $V_{out}=\qty{3.3}{\volt}$,
\qty{5}{\volt} and \qty{12}{\volt}.
\end{exercise}

%=======================================================================
\chapter{Switching regulators}
%=======================================================================

\begin{exercise}[Buck duty cycle]
A buck converter makes $V_{out}=\qty{5}{\volt}$ from
$V_{in}=(9\pm1)\,\unit{\volt}$. Give the nominal duty cycle $\delta$ and its range
over the input range.
\end{exercise}

\begin{exercise}[Boost from a battery]
A \qty{12}{\volt} rail must come from a battery $V_{in}=(5\pm0.5)\,\unit{\volt}$,
$f_{SW}=\qty{200}{\kilo\hertz}$.
(a)~Topology and the $\delta$ range.
(b)~With a capacitor-current swing $\Delta I_C=\qty{1.2}{\ampere}$ and
$\text{ESR}=\qty{50}{\milli\ohm}$, the output ripple.
(c)~Why is this topology's feedback loop deliberately made \emph{slow}?
\end{exercise}

\begin{exercise}[Minimum $L$ for CCM]
A buck has $V_{in}=\qty{9}{\volt}$, $V_{out}=\qty{5}{\volt}$ ($\delta=0.556$),
$f_{SW}=\qty{100}{\kilo\hertz}$, lightest load $I_{out,min}=\qty{50}{\milli\ampere}$.
Find the smallest $L$ that keeps it in CCM and give the next standard value.
\end{exercise}

%=======================================================================
\chapter{Laboratory: Buck converter}
%=======================================================================

\begin{exercise}[CCM or DCM?]
For a buck with the boundary (critical) current
$I_B=\dfrac{V_{out}(V_{in}-V_{out})}{2Lf_{sw}V_{in}}$, decide the conduction mode:
(a)~$V_{in}=\qty{8}{\volt}$, $V_{out}=\qty{5}{\volt}$, $R_L=\qty{30.5}{\ohm}$,
$I_B\approx\qty{0.11}{\ampere}$;
(b)~$V_{in}=\qty{24}{\volt}$, $V_{out}=\qty{5}{\volt}$, $R_L=\qty{50}{\ohm}$,
$I_B\approx\qty{0.23}{\ampere}$.
\end{exercise}

%=======================================================================
\chapter{Answers}\label{ch:answers}
%=======================================================================

Final results only --- the matching chapter of the \emph{manual} shows how each is
obtained.

\section*{Chapter 1 --- Specifications}
\ans{1.1}(a)~$S_I=0.010/1.495\approx\qty{6.7}{\milli\ohm}$;
$S_V=0.003/4=\qty{0.75}{\milli\volt\per\volt}$; ripple rejection
$=20\log_{10}(4/0.003)=\qty{62.5}{\dB}$.
(b)~$P_{out}=\qty{5}{\watt}$, $P_{in}\approx\qty{12}{\watt}$:
$\eta=\qty{41.7}{\percent}$, $P_d=\qty{7}{\watt}$ of heat.
\ans{1.2}Ripple rejection $=20\log_{10}(17.5/0.007)=\qty{68}{\dB}$;
$S_I=0.030/1.5=\qty{20}{\milli\ohm}$.

\section*{Chapter 2 --- The classical (transformer) power supply}
\ans{2.1}(a)~$C=I_{out}(T/2)/\Delta V=\qty{5}{\milli\farad}$.
(b)~$V_{MAX}=\qty{15}{\volt}$, secondary peak \qty{17}{\volt},
$V_{S,RMS}=\qty{12.0}{\volt}$; at $-\qty{10}{\percent}$ that is \qty{13.4}{\volt}
$\Rightarrow$ pick a \qty{15}{\volt_{RMS}} secondary, ratio $\approx230/15=15.3$.
(c)~Half-wave discharges for the whole period $T$, twice as long, so the same
$\Delta V$ needs $C=\qty{10}{\milli\farad}$.
\ans{2.2}$\Delta V=\qty{0.9}{\volt}$, $t_c=\qty{1.44}{\milli\second}$,
$I_{peak}=\qty{2.77}{\ampere}$.
\ans{2.3}$\Delta V'=\qty{0.43}{\volt}$ (smaller), $t_c'=\qty{0.98}{\milli\second}$,
$I_{peak}'=\qty{4.1}{\ampere}$ (larger) --- lower ripple, higher current peak.

\section*{Chapter 3 --- Power factor and its correction}
\ans{3.1}(a)~$\omega L=\qty{4}{\ohm}$, $|Z|=\qty{5}{\ohm}$,
$|I_S|=\qty{96}{\ampere}$, $\varphi=53.1^\circ$, $S=\qty{46.08}{\kilo\VA}$,
$P=\qty{27.67}{\kilo\watt}$, $PF=0.6$.
(b)~$C=L/(\omega^2L^2+R^2)=\qty{424.5}{\micro\farad}$; source current drops to
$\qty{58}{\ampere}$, $PF=1$.
\ans{3.2}(a)~$|I_S|=V_S/(\omega L)=\qty{120}{\ampere}$, $\varphi=90^\circ$,
$PF=0$. (b)~$C=1/(\omega^2L)=\qty{663.8}{\micro\farad}$ (resonance), ideally
$I_S=0$.

\section*{Chapter 4 --- Linear voltage regulators}
\ans{4.1}(a)~$P_{reg}=\qty{0.6}{\watt}$, $\eta=\qty{62.5}{\percent}$; junction
rise $\qty{36}{K}\Rightarrow T_j\approx\qty{61}{\celsius}$ --- fine without a heat
sink. (b)~$V_{MIN}=V_{out}+V_{DO,min}=\qty{7}{\volt}$.
\ans{4.2}$P_{reg}=\qty{7}{\watt}$, $\eta=\qty{42}{\percent}$; bare rise
\qty{420}{K} (far past the \qty{125}{\celsius} limit), with \qty{10}{K/W} rise
\qty{70}{K} $\Rightarrow T_j\approx\qty{95}{\celsius}$ (acceptable).
\ans{4.3}$\eta_{MAX}=\qty{62.3}{\percent}$ (\qty{3.3}{\volt}),
\qty{71.4}{\percent} (\qty{5}{\volt}), \qty{85.7}{\percent} (\qty{12}{\volt}).

\section*{Chapter 5 --- Switching regulators}
\ans{5.1}$\delta=V_{out}/V_{in}=\qty{55.5}{\percent}$; range
$\qty{50}{\percent}$ ($V_{in}=\qty{10}{\volt}$) to $\qty{62.5}{\percent}$
($V_{in}=\qty{8}{\volt}$).
\ans{5.2}(a)~$V_{out}>V_{in}$: \textbf{Boost}, $\delta=1-V_{in}/V_{out}$, range
$\qty{62.5}{\percent}\dots\qty{54.2}{\percent}$ (nominal \qty{58.3}{\percent}).
(b)~$\Delta V_{out}=\Delta I_C\cdot\text{ESR}=\qty{60}{\milli\volt}$.
(c)~The boost is prone to instability, so its loop is given a low bandwidth ---
it reacts slowly but will not oscillate.
\ans{5.3}$L_{min}=\dfrac{(V_{in}-V_{out})\delta}{2I_{out,min}f_{SW}}
=\qty{222}{\micro\henry}$; pick the next standard \qty{330}{\micro\henry}
($\Delta I_L\approx\qty{67}{\milli\ampere}$, CCM holds).

\section*{Chapter 6 --- Laboratory: Buck converter}
\ans{6.1}(a)~$I_{out}=V_{out}/R_L\approx\qty{0.16}{\ampere}>I_B$ $\Rightarrow$
\textbf{CCM}. (b)~$I_{out}\approx\qty{0.10}{\ampere}<I_B$ $\Rightarrow$
\textbf{DCM} (light load at a high step-down ratio).

\end{document}
